\begin{table}[htbp]
\centering\small
\caption{Posthoc conditioning/data controls: balanced H0 fixed-bin yield change, percentage points, mean $\pm$ sample SD over three paired fits. Relabeling also changes class-balanced loss weights. The matched-range comparison holds the physical edges identical; adding rows can still change minibatches and row weights.}
\begin{tabular}{lrr}\toprule
Model / contrast & UMA & MACE\\\midrule
AR / H0 range versus fixed & $+6.48\pm1.01$ & $+6.11\pm0.98$\\
AR / New rows / fixed edges & $-1.56\pm0.88$ & $-1.10\pm1.33$\\
AR / New rows / matched range edges & $-1.36\pm0.39$ & $-2.05\pm0.84$\\
CVAE $D=2$ / H0 range versus fixed & $+7.06\pm1.09$ & $+6.80\pm0.91$\\
CVAE $D=2$ / New rows / fixed edges & $+0.41\pm0.78$ & $+0.09\pm0.46$\\
CVAE $D=2$ / New rows / matched range edges & $+1.85\pm1.33$ & $+2.00\pm1.39$\\
\bottomrule\end{tabular}
\end{table}


A fresh, explicitly post-hoc control shows that most of the first-round
range-label gain does not require new structures. Starting from the same
original checkpoint and allocating the same 2,000-update continuation,
relabeling H0 alone improves balanced UMA yield by $6.48\pm1.01$ points for
AR and $7.06\pm1.09$ for $D=2$. Independent MACE assessment gives
$6.11\pm0.98$ and $6.80\pm0.91$ points. At identical first-round range edges,
adding the acquired rows instead changes UMA yield by $-1.36\pm0.39$ points
for AR and $+1.85\pm1.33$ for $D=2$; MACE gives $-2.05\pm0.84$ and
$+2.00\pm1.39$ points. These are three-pair sample SDs.

The larger improvement is therefore available through conditioning and its
associated loss reweighting on existing data. The smaller added-row effect
depends on architecture. This does not negate the measured expansion of
upper-energy samples: fixed-bin fidelity and exploration of a wider observed
energy range answer different questions. Neither outcome establishes DFT
accuracy, physical stability, or global coverage.
